\chapter{Modèles linéaires et de référence régularisés}\label{ch:ml:linear-and-regularised-baselines}

Une équipe passe un trimestre sur des arbres boostés par gradient pour la coupe transversale de 500 actions. Hors échantillon, sur dix ans,
leur portefeuille long-short par déciles affiche un ratio de Sharpe de 1.90 après coûts. Une régression lasso sur les mêmes quarante
caractéristiques classées, dont la pénalité est choisie par validation croisée, est ajustée en une seconde et atteint 1.91. Les arbres prévoient
mieux (un \omterm{def:ml:why-financial-machine-learning-is-different:r2}{R² hors échantillon} de 0.47~\% contre 0.32~\%), et sur un portefeuille qui échange les extrêmes du
classement, l'écart disparaît. Ce chapitre construit le modèle que tous les modèles suivants doivent battre, et mesure quand et
pourquoi il est difficile à battre.

\section{Le modèle à battre}

\begin{definition}[Modèle de référence]\label{def:ml:linear-and-regularised-baselines:baseline}
Un \emph{modèle de référence}\index{modèle de référence} (baseline) est le modèle le plus simple et crédible d'une tâche, ajusté et évalué avec les
mêmes données, découpages, cible et rapport que le candidat~: toute affirmation qu'un modèle fonctionne est une affirmation qu'il bat la
référence d'une marge supérieure au bruit de la comparaison.
\end{definition}

La référence est une discipline, non une formalité. Elle repère les fuites du chapitre~3 (une référence qui obtient un score aussi bon que le
candidat sur une cible qui devrait être difficile signale un problème dans le pipeline), elle chiffre le prix de la complexité, et c'est le modèle qui tourne
en production pendant que le candidat est en cours de validation. \texttt{firm.mlbase} est le harnais par lequel passe chaque modèle en coupe transversale du
livre~: variables classées par date, cible centrée par date (chapitre~1), ajustement sur les mois d'entraînement,
et un seul rapport pour les mois de test~: le R² par rapport à zéro, le coefficient d'information de rang moyen (rank IC) avec son $t$ HAC
(livre~7, chapitre~6), et le ratio de Sharpe d'un portefeuille long-short par déciles à pondération égale, brut et net de 20 points de base
par unité échangée.

Le panel du chapitre est celui de \texttt{firm.mlsynth}, avec quarante caractéristiques, dont six portent un rendement espéré injecté
(la moitié de sa variance étant non linéaire), et un ajout qui compte~: chaque caractéristique porte aussi un rendement du facteur
de moyenne nulle et de volatilité mensuelle de 2~\%. Un portefeuille trié sur des caractéristiques porte alors un risque de facteurs, comme le font les véritables,
et son ratio de Sharpe rejoint la fourchette publiée~: les portefeuilles déciles à pondération égale de Gu, Kelly et Xiu atteignent 2.45
pour leur meilleur réseau et 0.83 pour un modèle linéaire à trois variables. Sans les rendements des facteurs, les mêmes modèles
afficheraient des ratios de Sharpe proches de 5.

\section{Régression régularisée sur des variables classées}

Les moindres carrés ordinaires, la régression ridge, le lasso et l'elastic net relèvent du livre~4 (chapitre~16). Ce que ce chapitre
ajoute, c'est la façon de choisir leur pénalité et l'aspect que prend ce choix sur des données de marché.

\begin{definition}[Chemin de régularisation]\label{def:ml:linear-and-regularised-baselines:path}
Le \emph{chemin de régularisation}\index{chemin de régularisation} d'un estimateur pénalisé est sa solution, ou son score de validation,
en fonction de la pénalité $\lambda$~; la pénalité est choisie au meilleur score de validation croisée purgée du chemin.
\end{definition}

\begin{omfigure}
\begin{tikzpicture}
\begin{groupplot}[group style={group size=2 by 1,horizontal sep=1.6cm},width=6.6cm,height=5.2cm,
  tick label style={font=\footnotesize},label style={font=\small},grid=major,grid style={gray!25},table/col sep=comma,
  ymin=-0.02,ymax=0.42]
\nextgroupplot[xlabel={$\log_{10}\lambda$ (ridge)},ylabel={$R^2$ validé croisé (\%)},xtick={2,3,4,5},
  legend style={font=\footnotesize,at={(0.5,-0.33)},anchor=north,draw=none,fill=none},legend columns=2]
\addplot[omDef,very thick,mark=*] table[x=log10alpha,y=ridge]{figdata/ml/04-linear-and-regularised-baselines/path_ridge.csv};
\addplot[omProp,very thick,mark=triangle*] table[x=log10alpha,y=splines]{figdata/ml/04-linear-and-regularised-baselines/path_ridge.csv};
\legend{ridge,ridge sur splines}
\nextgroupplot[xlabel={$\log_{10}\lambda$ (lasso)},xtick={-4,-3.5,-3,-2.5},
  xticklabel style={/pgf/number format/fixed,/pgf/number format/precision=1}]
\addplot[omThm,very thick,mark=square*] table[x=log10alpha,y=lasso]{figdata/ml/04-linear-and-regularised-baselines/path_lasso.csv};
\end{groupplot}
\end{tikzpicture}

\omcaption{\omterm{def:ml:linear-and-regularised-baselines:path}{Chemins de régularisation}~: R² agrégé hors repli de trois replis purgés des vingt années d'entraînement, en
fonction de la pénalité. La régression ridge culmine à $\lambda = 10^4$ (replis d'entraînement d'environ 80\,000 lignes), le lasso à $10^{-3}$~; un cran plus loin, le
lasso annule chaque coefficient et score zéro. Données~: \texttt{ml\_baselines.paths}.}\label{fig:ml:baselines:path}
\end{omfigure}

Les chemins sont plats près de leurs sommets et chutent fortement de part et d'autre (\cref{fig:ml:baselines:path}). Au sommet, le
lasso conserve six coefficients~; un cran plus loin, il n'en conserve aucun, et un modèle qui prédit zéro partout obtient zéro,
ce qui vaut mieux que chacun des ajustements non pénalisés situés à gauche de la figure. Les moindres carrés sur 120\,000
lignes et quarante variables semblent un problème bien déterminé~; or, les rendements des facteurs faisant des mois, et non des lignes, l'échantillon
effectif (chapitre~1), il ne l'est pas, et la pénalité gagnante est grande.

\begin{method}[Choisir une pénalité sur un panel]\label{met:ml:baselines:penalty}
\begin{enumerate}
\item Classer chaque variable par date~; centrer la cible par date.
\item Construire des replis purgés de dates entières (les lignes d'une même date ne chevauchent jamais deux replis), avec un embargo égal à la longueur de l'étiquette.
\item Évaluer une grille logarithmique assez large pour que ses deux extrémités soient moins bonnes que le milieu~; si la meilleure valeur est à une extrémité,
élargir la grille.
\item Préférer la valeur la plus pénalisée à moins d'une erreur type du meilleur score (le plateau du chemin, et non son sommet).
\end{enumerate}
\end{method}

\section{Réduction de dimension~: composantes principales et moindres carrés partiels}

\begin{definition}[Régression sur composantes principales, moindres carrés partiels]\label{def:ml:linear-and-regularised-baselines:pcr}
La \emph{régression sur composantes principales}\index{régression sur composantes principales} (PCR) régresse la cible sur les
$k$ premières composantes principales des variables (livre~4, chapitre~22). Les \emph{moindres carrés partiels}\index{moindres carrés partiels}
(PLS) la régressent sur $k$ composantes choisies pour leur covariance avec la cible~: le premier vecteur de poids est
proportionnel à $X^\top y$, et chacun des suivants est trouvé de la même façon après avoir déflaté $X$ par les composantes déjà
retenues.
\end{definition}

\begin{proposition}[Ce que chaque réduction conserve]\label{prop:ml:baselines:pls}
Parmi les vecteurs unitaires $w$, $w\propto X^\top y$ maximise la covariance empirique de $Xw$ avec $y$, alors que la première direction principale
maximise la variance de $Xw$ quelle que soit $y$. Lorsque les variables sont non corrélées et de même variance, toutes les directions
ont la même variance et les composantes principales sont arbitraires~; la première composante du PLS est la direction de l'ajustement
par moindres carrés, à un facteur d'échelle près.
\end{proposition}

\begin{proof}
$\widehat{\Cov}(Xw, y)\propto w^\top X^\top y \le \lVert w\rVert\,\lVert X^\top y\rVert$ par Cauchy--Schwarz, avec égalité
pour $w\propto X^\top y$. Si $X^\top X = cI$, la variance empirique $w^\top X^\top Xw/n = c/n$ est la même pour tout $w$ unitaire~;
le coefficient des moindres carrés vaut $(X^\top X)^{-1}X^\top y = X^\top y/c$.
\end{proof}

Les rangs transversaux sont proches de ce cas~: quarante caractéristiques de même variance et peu corrélées. La PCR a besoin
de dix composantes avant de retrouver le signal et ne score encore que 0.13~\%~; la composante unique du PLS fait déjà ce que font
les moindres carrés, et les deux scorent de manière identique (\cref{tab:ml:baselines:table}). Sur de vraies caractéristiques, qui sont
corrélées par blocs (les mesures de valeur entre elles, les mesures de momentum entre elles), la PCR fait mieux, et Gu,
Kelly et Xiu ont obtenu 0.26~\% pour la PCR et 0.27~\% pour le PLS, au-dessus des 0.11~\% de leur elastic net.

\section{Modèles de référence pour la classification}

\begin{definition}[Régression logistique]\label{def:ml:linear-and-regularised-baselines:logistic}
Une \emph{régression logistique}\index{régression logistique} modélise $\P(y = 1\mid x) = 1/(1 + e^{-(b + x^\top\beta)})$ et estime
$(b,\beta)$ par maximum de vraisemblance, en général avec une pénalité ridge ou lasso sur $\beta$.
\end{definition}

Les étiquettes qui sont des signes (chapitre~2) appellent un classifieur~; sa référence naturelle est la \omterm{def:ml:linear-and-regularised-baselines:logistic}{régression logistique}. Sur le panel, elle
classe les actions aussi bien que les régressions (IC 0.060, comme pour la ridge), avec une probabilité en plus~: jeter la
taille de chaque rendement ne coûte presque rien lorsque cette taille est surtout du bruit.

\begin{definition}[Expansion sur une base de fonctions]\label{def:ml:linear-and-regularised-baselines:basis}
Une \emph{expansion sur une base de fonctions}\index{expansion sur une base de fonctions} remplace chaque variable $x_j$ par des fonctions $h_{j1}(x_j),\dots,h_{jm}(x_j)$
(splines, indicatrices de classes, polynômes) et ajuste un modèle linéaire sur celles-ci~; le modèle est additif en les variables et
non linéaire en chacune d'elles.
\end{definition}

La régression ridge sur des splines quadratiques, avec quatre nœuds par caractéristique, est la référence linéaire la plus forte du panel
(R² 0.35~\%)~: elle capte le carré et la valeur absolue que contient la vérité, mais non l'interaction, que
nul modèle additif ne peut capter.

\begin{table}[htbp]
\centering\small
\begin{tabular}{@{}lcccccc@{}}
\toprule
modèle (paramètre) & $R^2_{\mathrm{oos}}$ & IC & $t$ (HAC) & Sharpe brut & Sharpe net & rotation\\
\midrule
moindres carrés ordinaires & 0.25~\% & 0.060 & 5.6 & 1.85 & 1.61 & 0.83\\
ridge ($\lambda = 10^4$) & 0.30~\% & 0.060 & 5.6 & 1.83 & 1.59 & 0.83\\
lasso ($\lambda = 10^{-3}$) & 0.32~\% & 0.066 & 6.7 & 2.17 & 1.91 & 0.80\\
elastic net & 0.32~\% & 0.066 & 6.7 & 2.17 & 1.91 & 0.80\\
PCR (10 composantes) & 0.13~\% & 0.043 & 4.1 & 1.23 & 0.99 & 0.77\\
PLS (1 composante) & 0.25~\% & 0.061 & 5.6 & 1.87 & 1.62 & 0.82\\
\omterm{def:ml:linear-and-regularised-baselines:logistic}{régression logistique} & 0.22~\% & 0.060 & 5.7 & 1.84 & 1.59 & 0.84\\
ridge sur splines & 0.35~\% & 0.064 & 6.2 & 1.98 & 1.69 & 0.95\\
arbres boostés & 0.47~\% & 0.076 & 7.8 & 2.20 & 1.90 & 0.97\\
\midrule
plafond & 0.71~\% & & & & & \\
\bottomrule
\end{tabular}
\caption{Dix années de test du panel de quarante caractéristiques avec risque factoriel (500 actions, vingt années d'apprentissage). Ratios de
Sharpe du portefeuille de déciles extrêmes pondéré également, annualisés, nets de 20 points de base par unité échangée~;
la rotation est la part du carnet brut échangée chaque mois. Données~: \texttt{ml\_baselines.table}.}\label{tab:ml:baselines:table}
\end{table}

\section{Pourquoi la référence l'emporte souvent}

Le \Cref{tab:ml:baselines:table} classe les modèles selon le R², et le classement selon le ratio de Sharpe net n'est pas le même. Les
arbres boostés prévoient le mieux (0.47~\% contre 0.32~\% pour le lasso), mais leurs prévisions changent davantage d'un mois à l'autre
(rotation 0.97 contre 0.80), leur portefeuille supporte davantage de coûts, et aux extrémités du classement, là où un portefeuille
de déciles échange, les deux modèles s'accordent surtout. Le ratio de Sharpe net du lasso, 1.91, égale celui des arbres, 1.90.

La comparaison dépend aussi de la quantité d'historique disponible (\cref{fig:ml:baselines:crossing}). Entraînés sur les deux dernières
années précédant le test, les arbres ont un R² de $-0.71~\%$, alors que leur portefeuille de déciles affiche encore un ratio de Sharpe de
1.11 contre 0.90 pour la ridge~: avec peu de données, les arbres classent les actions passablement mais dimensionnent mal leurs prévisions, ce qui fait la
différence entre l'IC et le R² (\cref{prop:ml:why:ceiling}). Ils dépassent la ridge en R² entre cinq et
dix ans d'historique. Sur des données réelles, l'argument comporte une seconde moitié~: Kelly, Malamud et Zhou montrent que des modèles
fortement régularisés, avec plus de paramètres que d'observations, peuvent mieux prévoir que de petits modèles~; «~simple~» n'est donc pas la
thèse, c'est «~régularisé, et comparé à conditions égales~».

\begin{omfigure}
\begin{tikzpicture}
\begin{groupplot}[group style={group size=2 by 1,horizontal sep=1.6cm},width=6.6cm,height=5.2cm,
  xlabel={\small fenêtre d'apprentissage (mois)},xtick={24,60,120,240},tick label style={font=\footnotesize},
  label style={font=\small},grid=major,grid style={gray!25},table/col sep=comma,xmin=10,xmax=254]
\nextgroupplot[ylabel={$R^2_{\mathrm{oos}}$ (\%)},ymin=-0.8,ymax=0.8,
  legend style={font=\footnotesize,at={(1.1,-0.33)},anchor=north,draw=none,fill=none},legend columns=3]
\addplot[omThm,very thick,mark=square*] table[x=months,y=ridge]{figdata/ml/04-linear-and-regularised-baselines/crossing.csv};
\addplot[omDef,very thick,mark=*] table[x=months,y=boosting]{figdata/ml/04-linear-and-regularised-baselines/crossing.csv};
\addplot[black!60,thick,dashed] table[x=months,y=ceiling]{figdata/ml/04-linear-and-regularised-baselines/crossing.csv};
\legend{ridge,arbres boostés,plafond}
\nextgroupplot[ylabel={ratio de Sharpe net},ymin=0.6,ymax=2.4]
\addplot[omThm,very thick,mark=square*] table[x=months,y=ridge_sr]{figdata/ml/04-linear-and-regularised-baselines/crossing.csv};
\addplot[omDef,very thick,mark=*] table[x=months,y=boosting_sr]{figdata/ml/04-linear-and-regularised-baselines/crossing.csv};
\end{groupplot}
\end{tikzpicture}

\omcaption{Ridge et arbres boostés entraînés sur les 2 à 20 dernières années précédant les mêmes dix années de test. À gauche, R²
contre zéro~; à droite, ratio de Sharpe net du portefeuille de déciles. Le classement des arbres est utile tôt~; leurs dimensionnements sont
pas. Données~: \texttt{ml\_baselines.crossing}.}\label{fig:ml:baselines:crossing}
\end{omfigure}

\section{Tutoriel~: le modèle à battre}\label{tut:ml:linear-and-regularised-baselines:tut}

\begin{tutorial}
\textbf{Objectif.} Ajuster les modèles linéaires de référence le long de leurs \omterm{def:ml:linear-and-regularised-baselines:path}{chemins de régularisation} grâce à \texttt{firm.mlbase}, les comparer
aux arbres boostés dans un même rapport, et faire varier la fenêtre d'apprentissage. \textbf{État final~:} \cref{tab:ml:baselines:table},
\cref{fig:ml:baselines:path,fig:ml:baselines:crossing}.
\begin{tutsteps}
\item \textbf{Le rapport}~: R² contre zéro, le rang de l'IC avec son $t$ HAC, et le portefeuille de déciles brut et net.
\omcode{code/firm/mlbase/firm_mlbase.py}{61}{81}{Le portefeuille de déciles extrêmes long-short et le rapport standard.}\label{lst:ml:baselines:report}
\item \textbf{Le chemin}~: plis purgés de mois entiers, \omterm{def:ml:why-financial-machine-learning-is-different:r2}{R² hors échantillon} agrégé sur les plis.
\omcode{code/firm/mlbase/firm_mlbase.py}{84}{95}{Chemin de régularisation validé croisé.}\label{lst:ml:baselines:path}
\item \textbf{Exécuter} \texttt{ml\_baselines.paths()}, \texttt{table()}, \texttt{crossing()} et \texttt{fig\_baselines.py}.
\end{tutsteps}
\textbf{Pour aller plus loin.} Mettre la volatilité des facteurs à zéro et observer chaque ratio de Sharpe doubler~; ajouter
l'interaction $c_0c_3$ comme variable à la ridge sur splines et voir quelle part de l'avantage des arbres elle supprime.
\end{tutorial}

\section{Construction~: le harnais de modèles}\label{bld:ml:linear-and-regularised-baselines:mlbase}

\begin{build}
\bfield{Rôle} Chaque modèle en coupe transversale est comparé aux mêmes références, sur les mêmes découpages, avec le même
rapport.
\bfield{Interface} \texttt{rank\_features(X)}, \texttt{xs\_target(r, months)},
\texttt{fit(model, X, r, months, target)}, \texttt{predict(model, X, months)}, \texttt{long\_short(pred, r, q)},
\texttt{report(pred, r, q, cost\_bp, periods)} renvoyant \texttt{r2}, \texttt{ic}, \texttt{ic\_t}, \texttt{ls\_sr}, \texttt{ls\_sr\_net},
\texttt{turnover}~; \texttt{cv\_path(make, params, X, r, months, n\_folds, embargo)}~;
\texttt{Harness(make, target).run(X, r, train, test)}.
\bfield{Règles} Tout régresseur de type scikit-learn s'y branche~; les plis sont des dates entières~; le rapport n'est calculé que sur les dates de test
seulement~; les coûts sont imputés sur le notionnel échangé.
\bfield{Tests d'acceptation} \texttt{code/firm/mlbase/tests/}~: rangs et cibles centrées~; pondérations des déciles de somme nulle
avec un brut égal à deux~; le harnais retrouve un signal injecté et rapporte un brut supérieur au net~; le chemin préfère une pénalité
modérée sur des données bruitées.
\bfield{Pour aller plus loin} Déciles pondérés par la valeur~; le rapport par classe de taille~; une ridge à variables aléatoires dans
l'esprit de Kelly, Malamud et Zhou comme référence supplémentaire.
\end{build}

\begin{omsources}
\item S.~Gu, B.~Kelly et D.~Xiu, «~Empirical asset pricing via machine learning~», \emph{Review of Financial Studies}
33(5), 2020.
\item I.~Welch et A.~Goyal, «~A comprehensive look at the empirical performance of equity premium prediction~»,
\emph{Review of Financial Studies} 21(4), 2008.
\item B.~Kelly, S.~Malamud et K.~Zhou, «~The virtue of complexity in return prediction~», \emph{Journal of Finance}
79(1), 2024.
\item S.~de Jong, «~SIMPLS: an alternative approach to partial least squares regression~», \emph{Chemometrics and
Intelligent Laboratory Systems} 18, 1993.
\item A.~E.~Hoerl et R.~W.~Kennard, «~Ridge regression: biased estimation for nonorthogonal problems~»,
\emph{Technometrics} 12(1), 1970.
\end{omsources}

\section{Exercices}

\begin{exercise}[$\star$]\label{exo:ml:linear-and-regularised-baselines:1}
Un portefeuille par déciles de 500 actions détient 50 positions longues et 50 courtes, pondérées également, d'un brut égal à 2. Dans un mois, 30 titres changent
de chaque côté. Quelle part du livre brut est échangée, et combien cela coûte-t-il à 20 points de base par unité échangée~?
\end{exercise}

\begin{exercise}[$\star$]\label{exo:ml:linear-and-regularised-baselines:2}
Un portefeuille long-short rapporte en moyenne 1.6~\% par mois, avec une volatilité mensuelle de 2.5~\%. Quel est son
ratio de Sharpe annualisé~? Et après un coût mensuel de 0.2~\%~?
\end{exercise}

\begin{exercise}[$\star$]\label{exo:ml:linear-and-regularised-baselines:3}
Avec deux variables de même variance dont la corrélation vaut 0.9, quelle direction prend la première composante de l'ACP, et
et quand est-elle la bonne direction pour une régression~?
\end{exercise}

\begin{exercise}[$\star\star$]\label{exo:ml:linear-and-regularised-baselines:4}
Pourquoi les arbres boostés et le lasso obtiennent-ils le même ratio de Sharpe net alors que le R² des arbres est
supérieur de moitié~?
\end{exercise}

\begin{exercise}[$\star\star$]\label{exo:ml:linear-and-regularised-baselines:5}
Entraînés sur deux ans, les arbres ont un R² de $-0.71~\%$ et un ratio de Sharpe par déciles de 1.11. Concilier les deux nombres
avec la \cref{prop:ml:why:ceiling}.
\end{exercise}

\begin{exercise}[$\star\star$]\label{exo:ml:linear-and-regularised-baselines:6}
\emph{Trouver la faille.} «~Le R² de notre modèle profond dépasse de 0.2 point l'indice de référence linéaire publié, donc il est meilleur.~»
L'indice de référence publié a été estimé sur une autre période et un autre univers.
\end{exercise}

\begin{exercise}[$\star\star\star$]\label{exo:ml:linear-and-regularised-baselines:7}
\emph{Programmation.} Relancer \texttt{ml\_baselines.table} avec la volatilité des facteurs fixée à zéro (\texttt{STYLE = 0}). Indiquer
les ratios de Sharpe nets de la régression ridge et des arbres, et expliquer pourquoi ils changent autant alors que le R² bouge à peine.
\end{exercise}

\begin{exercise}[$\star\star\star$]\label{exo:ml:linear-and-regularised-baselines:8}
Montrer que, pour un plan orthonormé, le lasso annule tout coefficient dont la valeur des moindres carrés est inférieure
à $\lambda$ en valeur absolue, et s'en servir pour expliquer la chute à zéro du chemin du lasso dans la \cref{fig:ml:baselines:path}.
\end{exercise}

\section{Problème~: le modèle à battre}

\begin{problem}[{Problème du week-end --- le travail d'un trimestre face à celui d'une seconde}]\label{pb:ml:linear-and-regularised-baselines:1}
Le panel du chapitre~: 500 actions, quarante caractéristiques, risque factoriel, vingt ans d'apprentissage et dix ans de test.

\medskip\noindent\textbf{Partie I --- Le panel.}
\begin{enumerate}
\item Quel est le plafond, et quelles caractéristiques portent du signal~?
\item Pourquoi a-t-on ajouté les rendements des facteurs, et que font-ils aux ratios de Sharpe~?
\item Quels ratios de Sharpe par déciles, en pondération égale, Gu, Kelly et Xiu ont-ils rapportés~?
\item Pourquoi les moindres carrés sur 120\,000 lignes ont-ils besoin d'une pénalité~?
\end{enumerate}

\medskip\noindent\textbf{Partie II --- La famille linéaire.}
\begin{enumerate}[resume]
\item Quelles pénalités choisissent la régression ridge et le lasso, et qu'apporte le lasso en plus~?
\item Pourquoi la \omterm{def:ml:linear-and-regularised-baselines:pcr}{régression sur composantes principales} a-t-elle besoin de dix composantes et obtient-elle quand même 0.13~\%~?
\item Pourquoi les \omterm{def:ml:linear-and-regularised-baselines:pcr}{moindres carrés partiels} à une composante obtiennent-ils le même score que les moindres carrés~?
\item Que capture l'expansion sur une base de splines, et que ne peut-elle pas capturer~?
\end{enumerate}

\medskip\noindent\textbf{Partie III --- Les arbres face au \omterm{def:ml:linear-and-regularised-baselines:baseline}{modèle de référence}.}
\begin{enumerate}[resume]
\item Donner le R² des arbres et celui du lasso, l'IC et le ratio de Sharpe net.
\item Pourquoi le classement par ratio de Sharpe diffère-t-il du classement par R²~?
\item Quelle durée d'historique les arbres doivent-ils avoir pour dépasser la régression ridge en R²~?
\item Que rapportent les arbres avec deux ans d'historique, et pourquoi~?
\end{enumerate}

\medskip\noindent\textbf{Partie IV --- Le verdict.}
\begin{enumerate}[resume]
\item Énoncer le \emph{résultat nommé}~: le R² et l'IC de la régression ridge comparés à ceux des arbres, ainsi que l'historique à partir duquel les arbres
dépassent la régression ridge.
\item Quel modèle faudrait-il faire tourner en production pendant la validation des arbres, et pourquoi~?
\item Quelle référence supplémentaire le résultat sur la vertu de la complexité suggérerait-il~?
\item Combien de mois de résultats en conditions réelles faudrait-il pour départager le lasso et les arbres sur le ratio de Sharpe net~?
\item Qu'est-ce qui ferait de la \omterm{def:ml:linear-and-regularised-baselines:pcr}{régression sur composantes principales} la bonne référence~?
\item Quelle colonne du rapport faudrait-il montrer en premier à un gérant de portefeuille~?
\item Que perd-on en convertissant les rendements en signes~?
\item En une phrase~: à quoi sert un \omterm{def:ml:linear-and-regularised-baselines:baseline}{modèle de référence}~?
\end{enumerate}
\end{problem}

\section{Questions d'entretien}

\begin{interviewq}[$\star$ \iqroles{chercheur, ingénieur ML}]\label{iq:ml:linear-and-regularised-baselines:1}
On vous remet un nouveau jeu de données de rendements en coupe transversale. Quel est le premier modèle que vous ajustez, et comment le mettez-vous en place~?
\end{interviewq}

\begin{interviewq}[$\star$ \iqroles{chercheur}]\label{iq:ml:linear-and-regularised-baselines:2}
Régression ridge ou lasso pour quarante caractéristiques d'actions rangées~? Pourquoi~?
\end{interviewq}

\begin{interviewq}[$\star\star$ \iqroles{chercheur}]\label{iq:ml:linear-and-regularised-baselines:3}
Quelle est la différence entre la \omterm{def:ml:linear-and-regularised-baselines:pcr}{régression sur composantes principales} et les \omterm{def:ml:linear-and-regularised-baselines:pcr}{moindres carrés partiels}, et quand cela importe-t-il~?
\end{interviewq}

\begin{interviewq}[$\star\star$ \iqroles{chercheur, trader}]\label{iq:ml:linear-and-regularised-baselines:4}
Un modèle améliore le R² de 50~\% mais pas le ratio de Sharpe net du portefeuille. Comment cela est-il possible~?
\end{interviewq}

\begin{interviewq}[$\star\star$ \iqroles{ingénieur ML}]\label{iq:ml:linear-and-regularised-baselines:5}
Comment choisissez-vous concrètement la pénalité d'une régression ridge sur un panel d'actions~?
\end{interviewq}

\begin{interviewq}[$\star\star\star$ \iqroles{chercheur}]\label{iq:ml:linear-and-regularised-baselines:6}
Dérivez la première direction des \omterm{def:ml:linear-and-regularised-baselines:pcr}{moindres carrés partiels}, et dites pourquoi elle peut battre la première composante principale.
\end{interviewq}
